% Appendix: Model Onboarding Under Budget Constraints
% Include from paper/sections/appendix.tex

\subsection{Model Onboarding Under Budget Constraints}
\label{appendix:model_onboarding}

Production LLM portfolios are not static.
Providers release new models monthly---Gemini~2.5 Flash launched five
months after Gemini~2.5 Pro, GPT-4o-mini arrived three months after
GPT-4o---and platform teams need to incorporate them without
disrupting cost compliance or quality.
This section evaluates whether a single API call,
\texttt{register\_model()}, suffices to onboard a new model into a
\emph{budget-constrained} router.

\subsubsection{Why Care?}

The standard industry workflow for integrating a new model into a
supervised router is:
\begin{enumerate}[nosep,leftmargin=*]
  \item Collect the new model's responses on a representative prompt
    set and label them (days to weeks of human or LLM-judge
    annotation).
  \item Retrain the routing policy on the expanded labeled dataset.
  \item Update hard-coded cost tables and budget thresholds.
  \item Re-deploy and monitor for regressions.
\end{enumerate}
This process is slow, manual, and error-prone: stale cost tables
are the leading cause of budget violations in production
routers~(internal post-mortems).
Critically, step~1 requires a \emph{pre-collected evaluation dataset}
for the new model before it can receive any live traffic---an
offline data-collection campaign that must be repeated for every
model addition.

BanditGPT replaces steps 1--4 with a single call that seeds the new
arm with a T-shirt prior and lets the online learner discover its
quality--cost profile from live traffic.
The bandit still requires a reward signal, but that signal comes
from the same evaluation mechanism (LLM judge, user ratings, or
implicit behavioural composites) already serving the existing
portfolio---no dedicated offline dataset is collected for the new
model.
When the exploration mechanism routes a prompt to the new arm, the
existing reward pipeline evaluates its response at zero marginal
data-collection cost, and the bandit updates its posterior
immediately.
This also avoids two risks inherent to offline evaluation:
\emph{distribution mismatch} (a curated eval set may not reflect the
live prompt distribution) and \emph{dataset staleness} (the eval ages
as deployment patterns drift).
The BudgetPacer automatically accommodates the new arm's cost
without any reconfiguration.

\subsubsection{Experimental Design}

The same K=3 portfolio, data splits, and warmup priors used in
Experiments~01--03 are deployed with the BudgetPacer.
After Phase~1 online learning on the validation split ($n{=}1{,}785$),
\texttt{register\_model()} adds Gemini-2.5-Flash as a fourth arm
with \texttt{speed="fast"} and its published pricing
(input: \$0.30/M, output: \$2.50/M, blended: \$1.40/M).
Phase~2 continues online on the K=4 test split ($n{=}1{,}824$),
where Flash rewards are judged by DeepSeek-R1 using the same v3
continuous rubric as all other arms.

All conditions use the simulation-tuned hyperparameters from
Experiments~02--03: $\alpha{=}0.1$, $n_{\mathrm{eff}}{=}10$,
$\gamma{=}0.995$, consistent across all main-body experiments.

Two Phase~2 strategies are compared per budget tier:
\begin{itemize}[nosep,leftmargin=*]
  \item \textbf{Fixed Policy (uniform 1/4):}
    Equal allocation across all K=4 arms.
    The simplest possible onboarding strategy---no routing
    intelligence, no budget awareness.
  \item \textbf{BanditGPT (transfer):}
    Phase~1 posteriors carry over; Flash gets a T-shirt prior.
    BudgetPacer continues without reconfiguration.
\end{itemize}

\noindent
Three budget targets span the constraint regime:
tight ($B{=}\$2.3{\times}10^{-4}$/req),
moderate ($B{=}\$6.6{\times}10^{-4}$/req),
and loose ($B{=}\$1.9{\times}10^{-3}$/req).
An unconstrained BanditGPT baseline (no pacer) serves as a quality
ceiling.  Results: 20~seeds (9000--9019).

\subsubsection{Results (Figure~\ref{fig:model_onboarding})}

\begin{center}
\small
\begin{tabular}{@{}llrrrl@{}}
\toprule
\textbf{Strategy} & \textbf{Budget} &
  \textbf{Reward} & \textbf{Cost (\$/req)} &
  \textbf{Flash \%} & \textbf{Compliant?} \\
\midrule
Fixed Policy & any
  & 0.893 & \$4.05\text{e-}3 & 25.1\% & No (6--17$\times$) \\
\midrule
BanditGPT & tight
  & 0.853 & \$2.34\text{e-}4 & 10.7\% & \checkmark \\
BanditGPT & moderate
  & 0.889 & \$6.60\text{e-}4 & 19.4\% & \checkmark \\
BanditGPT & loose
  & 0.913 & \$1.80\text{e-}3 & 28.1\% & \checkmark \\
BanditGPT & unconstrained
  & 0.924 & \$7.68\text{e-}3 & 20.8\% & --- \\
\bottomrule
\end{tabular}
\end{center}

\paragraph{Flash is discovered in every seed.}
Across all 80 BanditGPT trials (20~seeds $\times$ 4~budget tiers),
Flash is first selected within ${\sim}340$ post-onboarding steps and
reaches sustained adoption (windowed share $>10$\% for 50
consecutive steps) by step~${\sim}370$.
Under the unconstrained ceiling, Flash stabilises at
$21.4 \pm 1.4$\% of traffic---comparable to its uniform-1/4 share
but \emph{contextually allocated} (the bandit routes Flash to
prompts where it matches Gemini-Pro quality at 10$\times$ lower
cost).

\paragraph{Budget determines adoption speed, not whether adoption occurs.}
At loose budgets, Flash reaches 33\% of traffic in the final 200
steps, comparable to the unconstrained ceiling (21\%).
At tight budgets, the pacer limits Flash to 10\%, just above the
sustained-adoption threshold---the router correctly identifies
that Flash (blended \$1.40/M) is cheaper than Mistral (\$1.00/M)
but more expensive than Llama (\$0.10/M), and allocates accordingly.

\paragraph{Flash displaces the right arm.}
Displacement analysis (early~200 vs.\ late~200 steps of Phase~2)
reveals a clear pattern:
\begin{itemize}[nosep,leftmargin=*]
  \item \textbf{Under budget pressure} (tight and moderate):
    Flash takes share primarily from \textbf{Mistral-Large}
    ($-15.1$\,pp at tight, $-15.5$\,pp at moderate).
    This is economically rational: Flash offers Mistral-level quality
    (${\sim}0.92$) at a lower blended cost (\$1.40 vs.\ \$1.00/M).
  \item \textbf{Unconstrained:} Flash takes share from
    \textbf{Gemini-Pro} ($-20.4$\,pp).  Without cost pressure, the
    bandit uses Flash as a near-Pro-quality alternative at
    $10\times$ lower cost.
\end{itemize}

\paragraph{The BudgetPacer requires no reconfiguration.}
Figure~\ref{fig:model_onboarding}(c) shows cost compliance
throughout the trajectory.
BanditGPT meets each budget target to within 1\% across both
phases, despite the portfolio expanding from K=3 to K=4
mid-trajectory.
The dual variable $\lambda_t$ adjusts smoothly: when Flash's
lower cost reduces the running average spend, $\lambda_t$ relaxes,
freeing budget for contextual exploration.

\paragraph{Fixed Policy: budget-oblivious baseline.}
The uniform-1/4 baseline highlights the cost of ignoring budget
constraints.
By giving Gemini-Pro 25\% of traffic, it spends \$4.05e-3 per
request---exceeding the tight budget by 17$\times$ and even the
loose budget by 2$\times$.
It achieves a reward of 0.893, which BanditGPT matches or exceeds
at moderate and loose budgets (0.889 and 0.913) while spending
6$\times$ and 2$\times$ less.

\subsubsection{Practical Implications}

\paragraph{One API call, zero downtime.}
At moderate and loose budgets---the regime of most production
deployments---onboarding a new model into BanditGPT requires a
single call to \texttt{register\_model()} with the model's pricing.
No offline evaluation, no policy retraining, no budget
reconfiguration.  The bandit discovers the model's quality--cost
niche from live traffic and the pacer auto-adjusts.

\paragraph{Budget-quality Pareto improvement.}
Flash is a Pareto-improving arm: it offers near-Pro quality at
Mistral-level cost.  A static router would require human analysis
to discover this; BanditGPT identifies and exploits it
automatically, shifting the portfolio's cost--quality frontier.

\paragraph{Tight budgets require deliberate exploration.}
Under tight budget constraints, the pacer's exploration budget is
limited, slowing (but not preventing) Flash adoption.
Practitioners running at very tight margins who want to onboard an
expensive newcomer should temporarily increase the budget target to
allow initial exploration, then tighten once the bandit calibrates.
